% sec-04: role authority.  Figure 2 is a plantuml-type use case diagram; the
% point is that every clinical authority stays with UC San Diego clinicians.
\section{Where Authority Would Sit}

The single question a site asks first is who decides. In this concept every
clinical authority sits with UC San Diego clinicians, and the applicant holds
none of it.

\begin{appfloat}[!tb]
\begin{appfig}[plantuml-type / role authority]
% Three actors left, one right; four use cases in a bounded system box.
% Ellipses 1.3 apart so association lines cannot cross one.
\node[umlpkg,minimum width=6.4cm,minimum height=6.6cm] (sys) at (5.0,-1.95) {};
\node[umlpkgtab,anchor=north west] at ([xshift=1mm,yshift=-1mm]sys.north west) {Clinical authority};
\node[umlusekey] (u1) at (5.0,-0.15) {Approve every robotic motion};
\node[umlusekey] (u2) at (5.0,-1.45) {Order conversion or abort};
\node[umlusecase] (u3) at (5.0,-2.75) {Decide dose escalation and holds};
\node[umlusecase] (u4) at (5.0,-4.05) {Adjudicate adverse-event attribution};
\umlactor{a1}{0}{-0.15}{Site PI, pancreatic surgery}
\umlactor{a2}{0}{-2.6}{GI medical oncologist}
\umlactor{a3}{0}{-4.6}{Clinical Trials Office and IRB}
\umlactor{a4}{10.2}{-3.2}{Applicant, ChemicalQDevice}
\draw[umlassoc] (a1) -- (u1.west);
\draw[umlassoc] (a1) -- (u2.west);
\draw[umlassoc] (a2) -- (u3.west);
\draw[umlassoc] (a2) -- (u4.west);
\draw[umlassoc] (a3) -- node[umlguard,pos=0.55] {oversight} (u4.west);
\node[umlnote,anchor=north west,text width=40mm] (n1) at (8.5,-0.4)
  {The applicant supplies the protocol, the software, the logging schema, and
  the analysis code, and holds no clinical decision right.};
\draw[umldash] (n1.west) -- (u1.east);
\pxmark{7.6}{-3.2}
\node[font=\tiny\sffamily,text=protogray,anchor=north,text width=30mm,align=center]
  at (7.6,-3.55) {no association with any clinical action};
\end{appfig}
\vspace{-0.7cm}
\figcaption{Where clinical authority sits. Every authorized action is associated
with a UC San Diego clinician, and the struck link shows the applicant has no
association with any of them.}
\end{appfloat}
